\documentclass[11pt,a4paper]{article}

\usepackage[utf8]{inputenc}
\usepackage[T1]{fontenc}
\usepackage[brazilian]{babel}
\usepackage[margin=2cm]{geometry}
\usepackage{enumitem}
\usepackage{hyperref}
\usepackage{titlesec}
\usepackage{parskip}

\hypersetup{
  colorlinks=true,
  linkcolor=black,
  urlcolor=black,
  pdftitle={Kamille Hartmann Maccagnan - Currículo},
  pdfauthor={Kamille Hartmann Maccagnan}
}

\titleformat{\section}{\large\bfseries}{}{0em}{}[\titlerule]
\titlespacing*{\section}{0pt}{12pt}{6pt}

\pagestyle{empty}
\setlength{\parindent}{0pt}
\setlength{\parskip}{6pt}

\newcommand{\cvheader}[4]{%
  \begin{center}
    {\LARGE\bfseries #1}\\[4pt]
    {\large #2}\\[6pt]
    #3\\[2pt]
    #4
  \end{center}
  \vspace{4pt}
}

\newcommand{\experience}[3]{%
  \textbf{#1} \hfill \textit{#2}\\[4pt]
  #3
  \vspace{8pt}
}

\begin{document}

\cvheader{Kamille Hartmann Maccagnan}%
{Analista de Data Analytics | Power BI | Indicadores e Dashboards | Excel}%
{Campo Comprido, Curitiba-PR \textbullet{} (41) 9 9928-4424 \textbullet{} \href{mailto:kamillehartmannm@gmail.com}{kamillehartmannm@gmail.com}}%
{\href{https://linkedin.com/in/kamille-hartmann-maccagnan/}{linkedin.com/in/kamille-hartmann-maccagnan}}

\section*{Resumo Profissional}

Profissional de Data Analytics e Business Intelligence com experiência na construção de dashboards em Power BI, desenvolvimento de indicadores de desempenho e elaboração de relatórios recorrentes para apoio à tomada de decisão. Atuação na análise de performance operacional, vendas e resultados, com foco em monitoramento de KPIs, identificação de inconsistências em bases de dados e padronização de informações para maior confiabilidade analítica. Experiência na estruturação de bases de dados em Excel, automação de fluxos de dados com Power Query e integração entre áreas de negócio, operações e tecnologia. Vivência em ambientes corporativos e multinacionais, com participação em reuniões de acompanhamento, documentação de processos e compartilhamento de análises com stakeholders. Perfil analítico e organizado, com prática em governança de dados, mapeamento de processos e capacitação de usuários para uso de ferramentas e indicadores.

\section*{Experiência Profissional}

\experience{Anjuss -- Analista de Suporte de TI}{07/2025 -- Atual}{%
Atuação no suporte a usuários e sistemas internos, com foco na organização de demandas e análise de dados para apoio à tomada de decisão. Responsável pela construção e manutenção de dashboards em Power BI para monitoramento de indicadores operacionais e desempenho de processos, além do desenvolvimento de relatórios recorrentes para acompanhamento de performance e visibilidade gerencial. Liderança do treinamento de novos funcionários e padronização de processos, com elaboração de documentações, POPs do sistema para as lojas e POPs técnicos para a área de T.I., garantindo consistência na mensuração e no uso das informações entre as áreas. Atuação na identificação de oportunidades de melhoria de processos e automação de rotinas analíticas, contribuindo para aumento de eficiência operacional, redução de retrabalho e disseminação do conhecimento entre as equipes.}

\experience{HORSE do Brasil -- Estagiária de TI (IS/IT LATAM)}{07/2024 -- 07/2025}{%
Atuação em múltiplas frentes corporativas, incluindo PMO, governança de dados, suporte e interface com áreas operacionais ligadas ao chão de fábrica, em ambiente multinacional. Acompanhamento de portfólio de projetos com foco em indicadores de desempenho, controle de cronogramas e suporte à tomada de decisão, com elaboração de relatórios executivos e dashboards para monitoramento de performance de projetos e iniciativas estratégicas. Atuação como interface entre áreas de negócio, fornecedores e stakeholders internacionais, participando de reuniões de acompanhamento conduzidas em até três idiomas simultaneamente e contribuindo para alinhamento de demandas, validação de entregas e compartilhamento de análises da operação. Participação em iniciativas de melhoria contínua, análise de riscos, documentação de processos e evolução de práticas de gestão e governança de dados.}

\experience{Restaurante Familiar -- Analista de Dados e Suporte Técnico}{01/2023 -- 07/2024}{%
Atuação na análise de vendas, performance operacional e resultados financeiros para suporte à gestão do negócio. Criação da base de dados do zero em Excel para estruturar informações e viabilizar o desenvolvimento de dashboards em Power BI, com definição de indicadores de desempenho para acompanhamento de vendas e resultados. Responsável pela padronização e organização de informações, contribuindo para maior confiabilidade dos dados e identificação de inconsistências que impactavam a análise. Realização de análises de performance e identificação de oportunidades de melhoria em processos operacionais, com atuação no alinhamento entre áreas operacionais e administrativas e apoio à tomada de decisão com base em indicadores-chave do negócio.}

\section*{Projetos e Conquistas}

Desenvolvimento de dashboards executivos para monitoramento de indicadores de desempenho e suporte à tomada de decisão. Estruturação de KPIs para análise de performance operacional e acompanhamento de resultados. Participação em iniciativas de melhoria contínua e otimização de processos com foco em eficiência operacional. Automação e padronização de fluxos de dados e relatórios em Excel e Power BI, reduzindo retrabalho e aumentando a confiabilidade das informações. 1º lugar no Transformation Day Renault 2023 com solução orientada a dados e melhoria de processos. Participação no Hackathon Bosch 2023 com foco em resolução de problemas de negócio com uso de dados.

\section*{Habilidades Técnicas}

\textbf{Data Analytics e Business Intelligence:} Power BI (DAX, Power Query), Excel, Google Sheets, SQL, análise de dados, visualização de dados, estruturação de bases de dados.

\textbf{Indicadores, Dashboards e Relatórios:} KPIs, dashboards executivos, relatórios gerenciais e recorrentes, monitoramento de performance, análise de performance operacional e financeira, tomada de decisão baseada em dados.

\textbf{Processos, Governança e Documentação:} Mapeamento e padronização de processos, melhoria contínua, governança de dados, documentação de processos, elaboração de manuais, POPs e materiais de apoio.

\textbf{Ferramentas:} Trello, Monday, ServiceNow, Spotfire, Pacote Office.

\textbf{Metodologias:} Scrum, Kanban, Boas Práticas PMOBOK.

\textbf{Atendimento e Capacitação:} Onboarding de usuários, treinamento e capacitação, suporte funcional, gestão de demandas, comunicação consultiva.

\textbf{Idiomas:} Inglês avançado \textbullet{} Espanhol básico \textbullet{} Coreano básico.

\section*{Formação Acadêmica}

\textbf{PUCPR} -- Graduação em Gestão da Tecnologia da Informação \hfill Previsão de conclusão: 06/2026\\
Curitiba, PR

\textbf{Escola Conquer} -- Pós-graduação em Liderança e Gestão em Tecnologia \hfill Conclusão: 2024\\
Curitiba, PR

\textbf{Universidade Tuiuti do Paraná} -- Graduação em Medicina Veterinária \hfill Conclusão: 06/2019\\
Curitiba, PR

\textbf{Qualittas} -- Pós-graduação em Clínica Médica e Cirúrgica de Animais Exóticos e Pets Não Convencionais \hfill Conclusão: 12/2022\\
Curitiba, PR

\end{document}
